% n-queens formalization report — Open Math Lab (mathforge), OPE-2.
% Build: ./build.sh  (tectonic, XeTeX engine; fonts from the tectonic bundle)
% Status: INTERNAL LAB REPORT. No external claim. Board (Paul) decides distribution.
\documentclass[11pt,a4paper]{article}

\usepackage{fontspec}
\usepackage{unicode-math}
\setmainfont{texgyrepagella}[Extension=.otf,UprightFont=*-regular,BoldFont=*-bold,ItalicFont=*-italic,BoldItalicFont=*-bolditalic]
\setmathfont{texgyrepagella-math.otf}
\setmonofont{DejaVuSansMono.ttf}[Scale=MatchLowercase]

\usepackage[a4paper,margin=27mm]{geometry}
\usepackage{amsthm}
\usepackage{booktabs}
\usepackage{longtable}
\usepackage{array}
\usepackage{fancyvrb}
\usepackage{xcolor}
\usepackage{enumitem}
\usepackage{microtype}
\usepackage[hidelinks]{hyperref}
\urlstyle{tt}
\emergencystretch=1.5em

\definecolor{codebg}{gray}{0.96}
\fvset{fontsize=\small,frame=single,framesep=2mm,rulecolor=\color{gray},xleftmargin=1mm}

\newtheorem{theorem}{Theorem}[section]
\newtheorem{lemma}[theorem]{Lemma}
\newtheorem{proposition}[theorem]{Proposition}
\theoremstyle{definition}
\newtheorem{definition}[theorem]{Definition}
\newtheorem{convention}[theorem]{Convention}
\theoremstyle{remark}
\newtheorem{remark}[theorem]{Remark}

\newcommand{\lean}[1]{\texttt{#1}}
\newcommand{\N}{\mathbb{N}}
\newcommand{\Fin}{\mathrm{Fin}\,}
\newcommand{\NonAttacking}{\mathrm{NonAttacking}}
\newcommand{\GoodEven}{\mathrm{GoodEven}}
\newcommand{\Repo}{\texttt{agentforce314/open-math-lab}}

\title{A Machine-Checked Proof of the \texorpdfstring{$n$}{n}-Queens Existence Theorem\\
\large A process report from Open Math Lab (mathforge)}
\author{Open Math Lab agent collective\thanks{Paperclip agents (Claude): Problem Scout, Formalist, Attack Lead, Adversarial Reviewer, Research Director (editor of this report). Human board: Paul. Repository: \url{https://github.com/agentforce314/open-math-lab}.}}
\date{12 September 2026 --- internal lab report, revision 1}

\begin{document}
\maketitle

\begin{abstract}
We report a zero-\lean{sorry} Lean~4 / Mathlib formalization of the classical
$n$-queens existence theorem: a placement of $n$ mutually non-attacking queens on an
$n\times n$ board exists if and only if $n\neq 2$ and $n\neq 3$
(Pauls 1874; explicit constructions by Hoffman--Loessi--Moore 1969 and
Bernhardsson 1991). The theorem is \emph{not new}; the contribution of this report is
(i)~the machine-checked artifact \lean{ProofLab.NQueensTheorem.n\_queens\_exists\_iff},
whose statement was frozen before any proof was attempted and whose axiom footprint is
exactly $\{\lean{propext}, \lean{Classical.choice}, \lean{Quot.sound}\}$;
(ii)~a 0-indexed transcription of Bernhardsson's construction in which every case
reduces to a linear-arithmetic goal closed by \lean{omega};
(iii)~an adversarial review with negative controls showing each clause of the
formal definition is load-bearing; and
(iv)~an account of the gate-by-gate process (five agent roles under a human board) that produced it in under one
hour of wall-clock time, including one transcription bug that survived a
brute-force check over $997$ boards and was caught only by the type checker.
The lab makes \textbf{no external claim}. Residual risks are listed explicitly.
\end{abstract}

\section*{Status and claim posture}
\addcontentsline{toc}{section}{Status and claim posture}
This is an internal report of Open Math Lab, produced under the lab's claim policy
(\path{docs/CLAIM_POLICY.md}): the default position is \emph{no claim}, and only the
human board may authorize external communication. Nothing here is a research result.
Every mathematical statement below has one of two labels: \textbf{Lean-checked}
(a named declaration in Appendix~\ref{app:lean}, verified by the kernel) or
\textbf{informal} (prose, evidence, or process narrative). The full process record,
including dead ends and an erratum, is preserved in the repository and summarized in
Section~\ref{sec:process}.

\tableofcontents

\section{Introduction}

``The $n$-queens problem'' names at least five different mathematical questions.
Table~\ref{tab:readings} lists the five that the lab's Problem Scout identified and
scored on the lab's feasibility rubric before any work was assigned. Exactly one of them,
the \emph{existence} question, is both settled in the literature and finitely encodable in
a way that a proof assistant can check in a single sprint. That one was funded; the others
are recorded as out of scope and remain untouched.

\begin{table}[ht]
\centering\small
\begin{tabular}{@{}l>{\raggedright\arraybackslash}p{7.2cm}r>{\raggedright\arraybackslash}p{4.2cm}@{}}
\toprule
& Reading & Score & Disposition \\
\midrule
(a) & Existence: a placement exists $\iff n\notin\{2,3\}$ & 88 & \textbf{target}; known-classical, formalize-only \\
(b) & Counting $Q(n)$ (OEIS A000170; $Q(8)=92$) & 54 & out of scope (no closed form known) \\
(c) & Asymptotics $Q(n)=((1\pm o(1))\,n e^{-\alpha})^n$ (Simkin 2021) & 28 & refused \\
(d) & Completion is NP-complete (Gent--Jefferson--Nightingale 2017) & 36 & refused \\
(e) & Toroidal queens exist $\iff \gcd(n,6)=1$ (Pólya 1918) & 80 & different theorem; out of scope \\
\bottomrule
\end{tabular}
\caption{Disambiguation by the Problem Scout (issue OPE-3). Scores are on the lab rubric
(five axes $\times$ 0--20). Only (a) was funded.}
\label{tab:readings}
\end{table}

The board's instruction was to ``solve the $n$-queens problem end to end using the whole
process, assuming the problem is new to you.'' We read ``new to you'' as an instruction
about \emph{process}, not about pretending: the repository already contained a small
Level~A artifact for boards $n\le 4$ (Section~\ref{sec:prior}), which every ticket cites as
prior art. What was genuinely new for the lab was the namesake theorem itself, which had
been explicitly left unproved.

\paragraph{What this report is.}
A formalization report and a process report. Section~\ref{sec:statement} fixes the formal
statement. Section~\ref{sec:construction} gives the mathematics of the construction in the
0-indexed form that was actually formalized, with each lemma tagged by its Lean name.
Section~\ref{sec:formalization} describes the Lean development and its verification.
Section~\ref{sec:review} reproduces the adversarial review, including the negative
controls. Section~\ref{sec:process} records the pipeline, gates, timeline, and budgets.
Section~\ref{sec:residuals} lists what is not claimed and what remains at risk.
Section~\ref{sec:lessons} draws process lessons. Appendix~\ref{app:lean} is the complete
Lean source; Appendices~\ref{app:transcripts}--\ref{app:timeline} are verification
transcripts and the ticket timeline.

\paragraph{What this report is not.}
It is not a claim of novelty, priority, or difficulty. It does not address readings
(b)--(e). It does not assert that the formal predicate $\NonAttacking$ is ``the''
definition of the problem; that identification is a human judgement, made independently
three times (Scout, Director, Reviewer) and exercised with counterexamples, but not
certified by the machine.

\section{Historical and formal context}
\label{sec:context}

\paragraph{History.}
Bezzel posed the eight-queens puzzle in 1848 and Nauck reported the $92$ solutions in
1850; the Gauss--Schumacher correspondence of that year concerns the $8\times 8$ count
only, not the general theorem. Pauls (1874) first proved existence for all $n\ge 4$.
Ahrens (1910) gave a textbook treatment. Hoffman, Loessi and Moore (1969) and
Bernhardsson (1991) gave explicit closed-form constructions split on $n \bmod 6$; the
latter's one-page note is the source transcribed here. Bell and Stevens (2009) survey the
field and confirm that (a) is settled, (b) is open-ended, and (e) is Pólya's theorem.

\paragraph{Mathlib.}
At the pinned Mathlib commit (\texttt{a719ba5c31}, Lean \texttt{v4.10.0}) a grep of
\texttt{Mathlib}, \texttt{Archive} and \texttt{Counterexamples} for the pattern
\begin{center}\small\texttt{nqueens|n\_queens|eightqueens|non.?attacking|queens?}\end{center}
returns no mathematical hit (the only match is a ``Queen Mary University'' comment). A negative control
(\lean{isTuranMaximal\_iff\_nonempty\_iso\_turanGraph}) fires, so the grep is sound.
This screen is local to the pin; it is not a survey of the wider Lean ecosystem.

\paragraph{Prior art in this repository.}
\label{sec:prior}
\texttt{ProofLab/NQueens.lean} (upstream PR~\#174) defines a \emph{list} encoding
\lean{IsNQueens : List (Fin n) → Prop} (length $n$, \lean{Nodup}, both diagonal families)
and proves \lean{queens\_two\_none}, \lean{queens\_three\_none}, and the two $n=4$
witnesses $[1,3,0,2]$ and $[2,0,3,1]$. The namesake theorem was left as a residual there.
The present work uses a \emph{function} encoding in a new file and re-proves the $n=2,3$
cases on that encoding, so no bridge lemma between the two encodings is needed---and none
was proved (Residual~1 in Section~\ref{sec:residuals}).

\section{The statement}
\label{sec:statement}

Throughout, boards are $0$-indexed: rows and columns are $\{0,\dots,n-1\}$, realised in
Lean as $\Fin\,n$. A \emph{placement} is a function $q:\Fin n\to\Fin n$; row $i$ holds
one queen, in column $q(i)$. One queen per row is thus built into the encoding.

\begin{definition}[\lean{NonAttacking}, Lean-checked definition]
\label{def:nonattacking}
A placement $q$ is \emph{non-attacking} if
\begin{enumerate}[nosep,label=(\roman*)]
\item (columns) $q$ is injective;
\item (anti-diagonals) for all $i\neq j$, $i+q(i)\neq j+q(j)$;
\item (main diagonals) for all $i\neq j$, $i+q(j)\neq j+q(i)$.
\end{enumerate}
All sums are in $\N$ after coercion from $\Fin n$; they are \emph{not} taken modulo $n$.
\end{definition}

Clause (iii) is the subtraction-free form of $i-q(i)\neq j-q(j)$: over $\mathbb{Z}$ the
two are the same inequality rearranged, and since every term is a natural number no
coercion subtlety arises. Writing it this way keeps the whole development inside
\lean{omega}'s fragment (linear arithmetic over $\N$ with literal coefficients).

\begin{convention}[$n=0$]
\label{conv:zero}
$\Fin 0\to\Fin 0$ has exactly one function and Definition~\ref{def:nonattacking} is
vacuous on it. The empty board therefore \emph{has} a placement. Consequently the theorem
is stated with right-hand side $n\neq 2\wedge n\neq 3$, which is total in $n$, rather
than the textbook ``$n=1$ or $n\ge 4$''. For $n\ge 1$ the two are the same theorem. This
convention was fixed by the Scout, confirmed by the Director, and re-derived independently
by the Reviewer.
\end{convention}

\begin{theorem}[\lean{n\_queens\_exists\_iff}, Lean-checked]
\label{thm:main}
For every $n\in\N$,
\[
\bigl(\exists\, q:\Fin n\to\Fin n,\ \NonAttacking(q)\bigr)\iff\bigl(n\neq 2\wedge n\neq 3\bigr).
\]
\end{theorem}

\begin{theorem}[\lean{queens\_exists\_of\_four\_le}, Lean-checked]
\label{thm:four}
For every $n\ge 4$ there is a non-attacking placement $q:\Fin n\to\Fin n$.
\end{theorem}

The Lean text of the definition and both theorem headers is reproduced verbatim below. It
was \emph{frozen} at the statement-pin gate (issue OPE-4) before any proof work began, and
was diffed byte-for-byte against the frozen text at the proof gate and again at review.

\begin{Verbatim}
def NonAttacking {n : ℕ} (q : Fin n → Fin n) : Prop :=
  Function.Injective q ∧
  ∀ i j : Fin n, i ≠ j →
    (i : ℕ) + (q i : ℕ) ≠ (j : ℕ) + (q j : ℕ) ∧
    (i : ℕ) + (q j : ℕ) ≠ (j : ℕ) + (q i : ℕ)

theorem queens_exists_of_four_le (n : ℕ) (hn : 4 ≤ n) :
    ∃ q : Fin n → Fin n, NonAttacking q

theorem n_queens_exists_iff (n : ℕ) :
    (∃ q : Fin n → Fin n, NonAttacking q) ↔ (n ≠ 2 ∧ n ≠ 3)
\end{Verbatim}

Theorem~\ref{thm:four} was split out at the pin gate as an optional target because it
isolates the only hard direction; the $n\le 3$ cases are finite.

\section{The construction}
\label{sec:construction}

This section is the mathematics as it was formalized. It follows the Attack Lead's log
(issue OPE-5), with the two corrections recorded in its erratum (Section~\ref{sec:erratum})
already applied. Every lemma is tagged with the Lean declaration that checks it. Throughout,
$h := \lfloor n/2\rfloor$; for even $n$, $n=2h$.

\subsection{Why three formulas rather than six residue classes}

Bernhardsson's construction is usually presented as a case split on $n \bmod 6$. The
Attack Lead chose a coarser partition: two closed formulas for even boards and one generic
\emph{corner extension} that handles all odd boards uniformly. The reason is tactical
rather than mathematical: every odd class then shares a single lemma whose only hypothesis
is that the even board has no queen on its main diagonal, and the whole proof stays in one
tactic flavour (\lean{split\_ifs <;> omega}). No modular arithmetic library (\lean{ZMod},
\lean{Nat.ModEq}) is needed anywhere. The alternative ``knight'' formula $r\mapsto 2r \bmod n$
for $\gcd(n,6)=1$ was machine-checked but deliberately not formalized, because it would
have required a second proof technique without removing any case from the plan.

\subsection{The even-board interface}

\begin{definition}[\lean{GoodEven}, Lean-checked definition]
For $n\in\N$ and $f:\N\to\N$, $\GoodEven(n,f)$ holds if, for all $r,s<n$:
\begin{enumerate}[nosep,label=(\alph*)]
\item $f(r)<n$;
\item $f(r)=f(s)\Rightarrow r=s$;
\item $r\neq s\Rightarrow r+f(r)\neq s+f(s)$;
\item $r\neq s\Rightarrow r+f(s)\neq s+f(r)$;
\item $f(r)\neq r$ \quad (\emph{no fixed point}: no queen on the main diagonal through the origin).
\end{enumerate}
\end{definition}

Clauses (a)--(d) say that $f$ restricted to $\{0,\dots,n-1\}$ is a non-attacking placement
in the sense of Definition~\ref{def:nonattacking}; clause (e) is the side condition the
corner extension needs.

\begin{lemma}[\lean{nonAttacking\_of\_nat}, Lean-checked; the only $\Fin$ lemma]
\label{lem:transfer}
If $f:\N\to\N$ satisfies (a)--(d) on $\{0,\dots,n-1\}$ then
$i\mapsto \langle f(i),\,\text{(a)}\rangle$ is a $\NonAttacking$ placement $\Fin n\to\Fin n$.
\end{lemma}

Everything else in the development is a statement about $\N\to\N$; Lemma~\ref{lem:transfer}
is the single crossing into dependent types.

\subsection{E1: \texorpdfstring{$n$ even, $n \equiv 0,4 \ (\mathrm{mod}\ 6)$}{n even, n = 0,4 (mod 6)}}

\begin{definition}[\lean{f₁}]
$f_1(r) = 2r+1$ if $r<h$, and $f_1(r)=2r-n$ otherwise.
\end{definition}

Odd columns $1,3,\dots,n-1$ fill the top half, even columns $0,2,\dots,n-2$ the bottom half.
For $n=4$ this is $[1,3,0,2]$, the Level~A witness.

\begin{lemma}[\lean{e1\_lt}, \lean{e1\_inj}, \lean{e1\_main},
\lean{e1\_no\_fix}; Lean-checked]
For even $n$, $f_1$ satisfies (a), (b), (d), (e).
\end{lemma}
\begin{proof}[Informal argument]
(a) $r<h\Rightarrow 2r+1\le n-1$; $r\ge h\Rightarrow 2r-n\le n-2$.
(b) Within a half, $2r+1=2s+1$ or $2r-n=2s-n$ forces $r=s$; across halves $2r+1=2s-n$
equates an odd and an even number.
(d) In $\mathbb{Z}$, $r-f_1(r)$ is $-r-1$ on the top half and $n-r$ on the bottom; within a
half this is injective, and across halves $-r-1=n-s$ gives $s-r=n+1>s$, impossible.
(e) $2r+1=r$ is impossible in $\N$; $2r-n=r$ gives $r=n$, contradicting $r<n$.
\end{proof}

\begin{lemma}[\lean{e1\_anti}, Lean-checked]
\label{lem:e1anti}
If $n\bmod 6\in\{0,4\}$ then $f_1$ satisfies (c).
\end{lemma}
\begin{proof}[Informal argument]
Within a half, $3r+1=3s+1$ or $3r-n=3s-n$ forces $r=s$. Across halves ($r<h\le s$),
$3r+1=3s-n$ gives $3(s-r)=n+1$, i.e.\ $n\equiv 2\pmod 3$, excluded by hypothesis.
\end{proof}

Lemma~\ref{lem:e1anti} is the only place the residue class is used, and the restriction is
sharp: for $n=8$, rows $1$ and $4$ of $f_1$ share anti-diagonal $4$.

\begin{proposition}[\lean{goodEven\_f₁}, Lean-checked]
If $n\bmod 6\in\{0,4\}$ then $\GoodEven(n,f_1)$.
\end{proposition}

\subsection{E2: \texorpdfstring{$n$ even, $n \equiv 2 \ (\mathrm{mod}\ 6)$}{n even, n = 2 (mod 6)}}

Here $h\bmod 3=1$, and for $n\ge 8$, $h\ge 4$. Bernhardsson's second formula contains a
``$x \bmod n$'' wrap; the Attack Lead resolved the wrap into explicit branches, since
$x\bmod n$ with variable $n$ lies outside \lean{omega}'s fragment.

\begin{definition}[\lean{f₂}]
\[
f_2(r)=
\begin{cases}
2r+h-1 & r<h,\ 2r\le h \quad\text{(P1)}\\
2r-h-1 & r<h,\ 2r> h \quad\text{(P2)}\\
2r+2-h & r\ge h,\ 2r<3h-2 \quad\text{(P3)}\\
2r+2-3h & r\ge h,\ 2r\ge 3h-2 \quad\text{(P4)}
\end{cases}
\]
\end{definition}

For $n=8$: $[3,5,7,1,6,0,2,4]$; for $n=14$:
$[6,8,10,12,0,2,4,9,11,13,1,3,5,7]$. The second half is the $180^\circ$ rotation of the
first, $f_2(n-1-r)=n-1-f_2(r)$, which is how the formula was derived but is not used by the
proof. Table~\ref{tab:e2} lists the linear forms that \lean{omega} sees in each piece.

\begin{table}[ht]
\centering\small
\begin{tabular}{@{}lllll@{}}
\toprule
piece & $f_2(r)$ & $r+f_2(r)$ (anti) & $r-f_2(r)$ (main, in $\mathbb Z$) & parity of $f_2(r)$ \\
\midrule
P1 & $2r+h-1$ & $3r+h-1$ & $-r-h+1$ & $h-1$ \\
P2 & $2r-h-1$ & $3r-h-1$ & $-r+h+1$ & $h-1$ \\
P3 & $2r-h+2$ & $3r-h+2$ & $-r+h-2$ & $h$ \\
P4 & $2r-3h+2$ & $3r-3h+2$ & $-r+3h-2$ & $h$ \\
\bottomrule
\end{tabular}
\caption{Piecewise-linear data for $f_2$. Values are as integers; the Lean definition
orders the operations so that no piece truncates in $\N$ on its own range
(Section~\ref{sec:erratum}).}
\label{tab:e2}
\end{table}

\begin{lemma}[\lean{e2\_lt}, \lean{e2\_inj}, \lean{e2\_anti}; Lean-checked]
If $n\bmod 6=2$ then $f_2$ satisfies (a), (b), (c).
\end{lemma}
\begin{proof}[Informal argument]
(a) Piece bounds: P1 $\le 2h-1=n-1$; P2 $\le h-3$; P3 $<n$; P4 $\le h$.
(b) Within a piece $2r+c=2s+c$; P1/P2 and P3/P4 cross-cases give $s-r=h$ against the
piece ranges; top-vs-bottom cross-cases have opposite parity (last column of
Table~\ref{tab:e2}).
(c) Within a piece $3r+c=3s+c$. Cross-piece constant differences are $2h$, $2h-3$,
$4h-3$, $3$, $2h-3$, $2h$; each $3(s-r)=\pm\text{const}$ forces $3\mid 2h$ or
$3\mid 4h$, impossible as $h\equiv 1\pmod 3$, except P2--P3 where $3(s-r)=-3$ gives
$s=r-1<h$, contradicting $s\ge h$.
\end{proof}

\begin{lemma}[\lean{e2\_main}, Lean-checked; needs $8\le n$]
If $n\bmod 6=2$ and $n\ge 8$ then $f_2$ satisfies (d).
\end{lemma}
\begin{proof}[Informal argument]
Within a piece, $-r+c=-s+c$. P1P2 and P3P4: $s-r=2h\ge n$. P1P4: $s-r=4h-3\ge n$.
P2P3: $r-s=3$ with $r<h\le s$. P1P3: $s-r=2h-3$ with $s<(3h-2)/2$ forces $h<4$.
P2P4: $s-r=2h-3$ with $r>h/2$, $s\le 2h-1$ forces $h\le 3$. Hence $h\ge 4$ suffices.
\end{proof}

\begin{lemma}[\lean{e2\_no\_fix}, Lean-checked; needs $4\le n$]
\label{lem:e2nofix}
If $n\bmod 6=2$ and $n\ge 4$ then $f_2$ satisfies (e).
\end{lemma}
\begin{proof}[Informal argument]
A fixed point would be P1 $r=1-h$; P2 $r=h+1>h$; P3 $r=h-2<h$; P4 $r=3h-2\ge 2h$
(using $h\ge 2$). Each contradicts its piece's range.
\end{proof}

The hypothesis $4\le n$ in Lemma~\ref{lem:e2nofix} is genuinely needed: at $n=2$,
$f_2(0)=0$ and $f_2(1)=1$. It was missing from the attack log's typed lemma
(Section~\ref{sec:erratum}); in the assembly it is discharged from $8\le n$.

\begin{proposition}[\lean{goodEven\_f₂}, \lean{goodEven\_exists}; Lean-checked]
If $n\bmod 6=2$ and $n\ge 8$ then $\GoodEven(n,f_2)$. Consequently every even $n\ge 4$
admits some $f$ with $\GoodEven(n,f)$.
\end{proposition}

(For even $n\ge 4$, $n\bmod 6\in\{0,2,4\}$; the class $2$ case has $n\ge 8$ automatically.)

\subsection{Odd boards: corner extension}

\begin{definition}[\lean{corner}]
For $m\in\N$ and $g:\N\to\N$, $\mathrm{corner}(m,g)(r)=g(r)$ if $r<m$, and $=m$ otherwise.
\end{definition}

Board $m+1$ is board $m$ with one more queen in the corner $(m,m)$.

\begin{lemma}[\lean{corner\_lt}, \lean{corner\_inj}, \lean{corner\_anti}, \lean{corner\_main}; all Lean-checked]
\label{lem:corner}
If $\GoodEven(m,g)$ then $\mathrm{corner}(m,g)$ satisfies (a)--(d) on the set
$\{0,\dots,m\}$.
\end{lemma}
\begin{proof}[Informal argument]
Both rows below $m$: inherited from $g$. One row equal to $m$:
(a) $m<m+1$; (b) $g(r)=m$ contradicts $g(r)<m$; (c) $r+g(r)\le 2m-2<2m=m+m$;
(d) $r+m\neq m+g(r)\iff g(r)\neq r$, which is exactly clause (e) of $\GoodEven$.
\end{proof}

The last line is the whole point of clause (e): without it the corner queen would lie on
the main diagonal through some $g$-queen. This is also why a toroidal (modular) solution
cannot be corner-extended (Section~\ref{sec:deadends}).

\subsection{Assembly}

\begin{proof}[Proof of Theorem~\ref{thm:four} (\lean{queens\_exists\_of\_four\_le})]
Write $n=2k$ or $n=2k+1$. If even, take the $f$ from \lean{goodEven\_exists} and apply
Lemma~\ref{lem:transfer}. If odd, $2k\ge 4$; take $g$ with $\GoodEven(2k,g)$, form
$\mathrm{corner}(2k,g)$, apply Lemma~\ref{lem:corner} and then Lemma~\ref{lem:transfer}.
\end{proof}

\begin{proof}[Proof of Theorem~\ref{thm:main} (\lean{n\_queens\_exists\_iff})]
($\Rightarrow$) $\Fin 2\to\Fin 2$ has $4$ functions and $\Fin 3\to\Fin 3$ has $27$; after
unfolding $\NonAttacking$, kernel \lean{decide} refutes each
(\lean{queens\_two\_none'}, \lean{queens\_three\_none'}). ($\Leftarrow$) If $n<4$ then
$n\in\{0,1\}$ and the identity is a placement (by \lean{decide}); otherwise
Theorem~\ref{thm:four}.
\end{proof}

\subsection{Dead ends (informal)}
\label{sec:deadends}

The attack log records three routes that were tried and rejected, with the obstruction in
each case:
\begin{enumerate}[nosep]
\item \emph{$1$-indexed ``list surgery''} for $n\bmod 6\in\{2,3\}$ (``swap $1\leftrightarrow 3$,
move $5$ to the end''). Transcribed to $0$-indexed form it is a five-piece function with
three single-point pieces, and the $n\bmod 6=3$ list is \emph{not} a corner extension of the
$n\bmod 6=2$ list (for $n=9$ the last row holds column $2$, not $8$). Rejected for
Bernhardsson's closed formula, whose companion \emph{is} the corner extension. Obstruction:
more cases and no reuse, not a mathematical failure.
\item \emph{Corner-extending a toroidal solution.} Any modular solution
$r\mapsto (ar+b)\bmod m$ has a queen on every wrapped diagonal, in particular on the actual
main diagonal $r=q(r)$, so clause (e) fails and the corner queen is attacked. Exact
obstruction.
\item \emph{A single $\bmod n$ formula for even $n$.} $(2r+c)\bmod n$ is never injective
for even $n$ ($\gcd(2,n)=2$), so no one-line modular argument exists for even boards; the
$r<h$ split is unavoidable.
\end{enumerate}

\section{The formalization}
\label{sec:formalization}

\subsection{Architecture}

The development is one file, \path{proofs/lean-project/ProofLab/NQueensTheorem.lean}
(272 lines, Appendix~\ref{app:lean}), in namespace \lean{ProofLab.NQueensTheorem}. It
contains five definitions ($\NonAttacking$, $\GoodEven$, $f_1$, $f_2$, $\mathrm{corner}$)
and twenty-two theorems. Its shape follows the attack plan exactly:
\begin{itemize}[nosep]
\item one transfer lemma from $\N\to\N$ into $\Fin n\to\Fin n$ (Lemma~\ref{lem:transfer});
\item an interface record $\GoodEven$ so that E1, E2 and the corner step compose without
knowing each other's formulas;
\item every arithmetic lemma proved by the same three-token script
\lean{simp only [f]; split\_ifs <;> omega}, one line each;
\item the two finite cases by kernel \lean{decide} on the function space, after
\lean{unfold NonAttacking} (plain \lean{decide} does not see through the \lean{def});
\item three \lean{example}s as sanity checks: the Level~A witness satisfies the predicate,
$f_1(4)=[1,3,0,2]$, $f_2(8)=[3,5,7,1,6,0,2,4]$.
\end{itemize}
The \lean{unusedVariables} linter is disabled for the arithmetic section because
\lean{omega} consumes the bound hypotheses $r<n$, $s<n$ from context without a syntactic
reference; the Formalist checked with \lean{clear} that they are load-bearing (e.g.\
\lean{e1\_main} fails at $n=4$, $r=0$, $s=5$ without $s<n$).

\subsection{Verification}
\label{sec:verification}

The proof was built, from a deleted module cache, by three parties on the same machine
(Lean 4.10.0, arm64 macOS; Mathlib \texttt{a719ba5c31}): the Director at the proof gate,
the Reviewer during review, and the Director again on the integrated close-out branch.
Each time:
\begin{Verbatim}
✔ [2240/2240] Built ProofLab.NQueensTheorem
Build completed successfully.
\end{Verbatim}
with no \lean{sorry} warning, and
\begin{Verbatim}
'ProofLab.NQueensTheorem.n_queens_exists_iff' depends on axioms:
    [propext, Classical.choice, Quot.sound]
'ProofLab.NQueensTheorem.queens_exists_of_four_le' depends on axioms:
    [propext, Quot.sound]
\end{Verbatim}
A grep of the file for the pattern
\begin{center}\small\texttt{sorry|admit|axiom|native\_decide|implemented\_by|extern|unsafe|partial}\end{center}
returns only the expected \lean{decide} calls. \lean{Classical.choice} enters through
Mathlib's \lean{Fintype} decidability instances used by \lean{decide} on the finite
function spaces; \lean{queens\_exists\_of\_four\_le} does not need it.

\subsection{Two deviations from the attack log, and an erratum}
\label{sec:erratum}

The Formalist reported two places where the landed Lean differs from the Attack Lead's
\emph{typed} lemma list; the Reviewer confirmed both independently.

\paragraph{(a) $\N$-truncation in the typed $f_2$.}
The log typed pieces P3/P4 as \lean{2 * r - n / 2 + 2} and \lean{2 * r - 3 * (n / 2) + 2}.
In $\N$ these parse as $(2r-3h)+2$ and truncate whenever $2r\in\{3h-2,3h-1\}$: for
$n=8,r=5$ the typed form gives $2$ instead of $0$; for $n=14,r=10$ it gives $2$ instead of
$1$. The Reviewer's own script found $166$ mismatched rows over $8\le n\le 1000$,
$n\equiv 2\pmod 6$---exactly one row per board---and the typed form would have broken
injectivity (at $n=8$, rows $5$ and $6$ would collide). The Attack Lead's Python check used
signed integers and so never exercised the $\N$ form. Lean's $f_2$ is written
\lean{2 * r + 2 - n / 2} and \lean{2 * r + 2 - 3 * (n / 2)}, which agree with the
signed formula on every branch where they are evaluated. The mathematics in the log and the
brute-force evidence were correct; only the $\N$ transcription was wrong. The log now
carries an erratum header and the corrected forms.

\paragraph{(b) A missing hypothesis.}
The typed \lean{e2\_no\_fix} had no size hypothesis and is false at $n=2$; the log's
justification ``(as $h\ge 2$)'' was unsupported. Lean adds $4\le n$, discharged by
\lean{omega} from the assembly's $8\le n$. The Reviewer tabulated every landed declaration
against the log (Table~\ref{tab:hypdiff}): this is the only hypothesis delta in the file,
and nothing was weakened.

\begin{table}[ht]
\centering\small
\begin{tabular}{@{}lll@{}}
\toprule
declaration(s) & hypotheses in log & delta in Lean \\
\midrule
\lean{nonAttacking\_of\_nat} & \lean{hlt hinj hanti hmain} & none \\
\lean{e1\_lt}, \lean{e1\_inj}, \lean{e1\_main}, \lean{e1\_no\_fix} & $n\bmod 2=0$ & none \\
\lean{e1\_anti}, \lean{goodEven\_f₁} & $n\bmod 6\in\{0,4\}$ & none \\
\lean{e2\_lt}, \lean{e2\_inj}, \lean{e2\_anti} & $n\bmod 6=2$ & none \\
\lean{e2\_main}, \lean{goodEven\_f₂} & $n\bmod 6=2$, $8\le n$ & none \\
\lean{e2\_no\_fix} & $n\bmod 6=2$ & \textbf{$+\,4\le n$} (needed) \\
\lean{goodEven\_exists} & $4\le n$, $n\bmod 2=0$ & none \\
\lean{corner\_*} & sketched & made concrete; none \\
\lean{queens\_exists\_of\_four\_le} & $4\le n$ (frozen) & none \\
\bottomrule
\end{tabular}
\caption{Hypothesis diff, attack log versus landed Lean (from the adversarial review).}
\label{tab:hypdiff}
\end{table}

\section{Adversarial review}
\label{sec:review}

The lab's rule is that no mathematical claim is closed without \emph{both} a green Lean
build \emph{and} written approval from an Adversarial Reviewer with veto power and a
non-empty residual-risk list. The Reviewer (issue OPE-7) was given the artifact pinned to a
commit and a binding checklist; the review file and its scratch Lean file are in the
repository (\path{attacks/n-queens-20260912-180329/ADVERSARIAL_REVIEW.md},
\path{review_checks.lean}). Verdict: \textbf{approve with residuals}.

\subsection{Statement fidelity}

The Reviewer read Definition~\ref{def:nonattacking} cold and matched each clause to the
chess rule (rows = domain, columns = codomain, injectivity = distinct columns, $i+q(i)$ =
anti-diagonal index, $i+q(j)=j+q(i)$ = shared main diagonal, sums in $\N$ not
$\Fin n$). Two properties were then established by kernel \lean{decide}:
\begin{itemize}[nosep]
\item \emph{Not vacuous:} the predicate is unsatisfiable at $n=2,3$.
\item \emph{Not over-strong:} concrete boards for $n=4,5,6,8,14$ are accepted, including
$[1,3,0,2]$, which would fail twice if the sums were taken modulo $4$.
\end{itemize}

\subsection{Negative controls}

To show that each clause is load-bearing, the Reviewer enumerated all $256$ functions
$\Fin 4\to\Fin 4$ (informally, in Python) to select boards that violate \emph{exactly one}
clause, then checked in Lean that $\NonAttacking$ rejects each, and that the predicate with
that one clause removed accepts it (Table~\ref{tab:controls}). The enumeration also confirms
that exactly two of the $256$ functions are accepted, the two classical $n=4$ solutions.

\begin{table}[ht]
\centering\small
\begin{tabular}{@{}llll@{}}
\toprule
board & violates only & $\NonAttacking$ & weakened predicate accepts \\
\midrule
$[0,1,2,3]$ & main diagonal & rejected & \lean{NoMain} \\
$[3,2,1,0]$ & anti-diagonal & rejected & --- \\
$[1,3,2,0]$ & anti-diagonal & rejected & \lean{NoAnti} \\
$[0,0,0,0]$ & injectivity & rejected & \lean{NoInj} \\
$[1,1,0,2]$ & injectivity \emph{and} anti-diagonal & rejected & (not single-clause; replaced) \\
\bottomrule
\end{tabular}
\caption{Negative controls, all by kernel \lean{decide} (\texttt{review\_checks.lean}).
Any control accepted by $\NonAttacking$ would have been an automatic reject.}
\label{tab:controls}
\end{table}

\subsection{Machine-trust boundary}

``Lean checked it'' means: the Lean 4.10.0 kernel accepted a term of type
\[\forall n,\ (\exists q,\ \NonAttacking\,q)\leftrightarrow(n\neq 2\wedge n\neq 3)\]
from the three standard axioms plus the pinned Mathlib's definitions. This certifies that the formal
proposition follows from those axioms, modulo trusting the kernel, the \lean{elan}-installed
toolchain binary, the Mathlib \texttt{.olean} cache on one machine, and the \lean{lake}
build graph. It does \emph{not} certify that $\NonAttacking$ is the right predicate---that
is the human judgement documented above---nor anything about readings (b)--(e), nor any
link to the Level~A list encoding. No \lean{native\_decide} or \lean{ofReduceBool}
shortcut was used anywhere.

\section{Process}
\label{sec:process}

\subsection{Pipeline and gates}

The lab runs as five Paperclip agent roles under a human board. For a board-named problem
the Scout shortlist gate is satisfied by the board, but every other gate still fires. Each
specialist was woken \emph{one at a time}; the Director re-verified the artifact at each gate
before releasing the next ticket, and the Director's own checks were logged on the parent
issue so the Reviewer could see what was and was not yet independent.

\begin{enumerate}[nosep]
\item \textbf{Board} (OPE-2): names the problem; asks for the whole process; approves
user-local tooling.
\item \textbf{Director}: installs Lean 4.10.0 via \lean{elan} (one macOS-26 linker
workaround, documented); frames the readings; opens the child chain; sets kill criteria.
\item \textbf{Problem Scout} (OPE-3): disambiguation, rubric score, literature handles,
Mathlib screen, Level~B statement pin, top three definition risks, budget. Does not solve.
\item \textbf{Formalist, statement pin} (OPE-4): the frozen text compiles with a single
\lean{sorry}; \lean{\#print axioms} shows \lean{sorryAx}, as it must for a stub.
\item \textbf{Attack Lead} (OPE-5): 0-indexed construction, informal proofs shaped for
\lean{omega}, brute-force evidence to $n=1000$, Lean-shaped lemma list, dead-end map.
Labelled \emph{status: informal}.
\item \textbf{Formalist, proof} (OPE-6): closes both theorems in one wave (cap: two), reports
the two deviations.
\item \textbf{Adversarial Reviewer} (OPE-7): fresh rebuild, axioms, fidelity, negative
controls, hypothesis diff, scope sweep, residuals. Veto power.
\item \textbf{Director close-out} (OPE-2): integrates branches, records the erratum, moves
the ledger and catalog to \emph{formalized}, prepares the claim packet
(\texttt{mathforge claim prepare}) with recommendation \emph{no claim}.
\end{enumerate}

\subsection{Timeline and budget}

All work took place on 12 September 2026. From the board's comment (17:37 UTC) to the
close-out disposition (18:34 UTC) is 57 minutes of wall-clock time across eight
heartbeats; the per-ticket commit times are in Appendix~\ref{app:timeline}. Token budgets
were set per ticket (Scout and stub small; Attack Lead cap 60k, reported $\approx$45k used;
Formalist proof wave 1 $\approx$120k of a two-wave cap; Reviewer $\approx$80k); the Scout's
overall estimate of $\approx$300k was accurate to within a few percent. The Director's kill
criterion (stop if more than one of the four proof blocks remained open after wave 2) was
never triggered.

Two facts keep this timeline in proportion. The theorem is classical and its construction
is a one-page note in the literature; and the repository already had the Lean toolchain
configuration, a Mathlib cache, a claim policy, a review checklist, and a role definition
for every participant. The hour measures the \emph{process}, not the mathematics.

\subsection{Artifacts}

All artifacts are on the fork \Repo, branch \texttt{ope/2-nqueens-closeout}
(pull request \#8, which integrates \#3, \#5, \#6, \#7):
\begin{itemize}[nosep]
\item \texttt{catalog/problems/n-queens/STATEMENT.md}, \texttt{DOSSIER.json} --- Scout dossier and pin;
\item \path{attacks/n-queens-20260912-180329/} --- attack log (with erratum), strategies,
\path{verify_construction.py} and its output, the adversarial review, \path{review_checks.lean}
and its output;
\item \path{proofs/lean-project/ProofLab/NQueensTheorem.lean} --- the proof;
\item \texttt{problems/n-queens/CLAIM\_PACKET.md} --- claim packet (no claim);
\item \texttt{docs/PROBLEM\_LEDGER.md} --- lifecycle record;
\item \texttt{proofs/n-queens/paper/} --- this report and its build script.
\end{itemize}

\section{What is not claimed; residual risks}
\label{sec:residuals}

\paragraph{Not claimed.}
New mathematics; priority in Mathlib or elsewhere; anything about readings (b)--(e);
that $\NonAttacking$ is the unique correct formalization; equivalence with the Level~A
list encoding; any ``AI solved $n$-queens'' framing.

\paragraph{Residual risks (from the review, board-facing).}
\begin{enumerate}[nosep]
\item No Lean bridge \lean{IsNQueens (List.ofFn q) ↔ NonAttacking q}; two encodings coexist.
\item Only reading (a) is proved.
\item Definition fidelity is a human judgement, made three times; a further reader is cheap.
\item Toolchain trust is single-machine; no \lean{lean4checker} re-check of the \texttt{.olean}.
\item No priority or novelty; the Mathlib screen is local to the pin.
\item The attack log's erratum (Section~\ref{sec:erratum}), now fixed, is kept on record.
\item $n\neq 2\wedge n\neq 3$ includes the vacuous $n=0$ (Convention~\ref{conv:zero}); do not
paraphrase as ``$n=1\vee n\ge 4$'' without the caveat.
\item An unused \lean{import ProofLab.NQueens} in the new file; harmless, since the axiom
print rather than the build is the gate.
\end{enumerate}

\section{Lessons}
\label{sec:lessons}

\begin{enumerate}[nosep]
\item \textbf{Brute-force evidence must run on the typed formula.} The Attack Lead's
check evaluated the pin's three clauses literally over all ordered pairs on $997$ boards and
passed on the first run---and still missed a bug, because it computed with signed integers
while the lemma list was typed in $\N$. Future attack tickets should evaluate the exact
$\N$-typed expression (or generate the Python from the Lean text).
\item \textbf{Freeze the statement before the proof.} Diffing the frozen text at every
gate cost seconds and removed an entire class of ``proved a different theorem'' failures.
\item \textbf{Isolate the hard direction.} Splitting off \lean{queens\_exists\_of\_four\_le}
at the pin gate let the finite cases, the construction, and the review proceed as separate
units of work.
\item \textbf{One interface, one tactic.} The $\GoodEven$ record and the
\lean{split\_ifs <;> omega} discipline meant that the Formalist's wave was almost entirely
transcription, and that the Reviewer could diff hypotheses mechanically.
\item \textbf{The review gate is not a formality.} Both deviations were real defects in a
document that had passed a machine check; they were caught by the type checker and
independently confirmed by the Reviewer, and the ticket that was funded to prevent this
class of error (the transcription step) was the one that introduced it.
\item \textbf{Kernel \lean{decide} is enough for small finite cases} ($4$ and $27$
functions) and keeps the axiom footprint standard; \lean{native\_decide} was never needed
and is not an accepted gate.
\end{enumerate}

\section{Conclusion}

The lab set out to run its full process on a problem chosen by the board, and did so: a
scouted and disambiguated target, a frozen statement, an informal construction with
machine evidence and an honest dead-end map, a zero-\lean{sorry} Lean proof with a standard
axiom footprint, an adversarial review with negative controls, and a close-out with a
claim packet recommending no claim. The theorem was known in 1874. What the lab can now
say, and could not before, is that the specific formal proposition in
Section~\ref{sec:statement} follows from Lean's three standard axioms, and that the
definition it is about has been read and stress-tested by three independent readers.

\begin{thebibliography}{9}
\bibitem{ahrens} W.~Ahrens, \emph{Mathematische Unterhaltungen und Spiele}, 2nd ed., Teubner, 1910.
\bibitem{bellstevens} J.~Bell and B.~Stevens, ``A survey of known results and research areas for $n$-queens,'' \emph{Discrete Mathematics} 309(1):1--31, 2009.
\bibitem{bernhardsson} B.~Bernhardsson, ``Explicit solutions to the $N$-queens problem for all $N$,'' \emph{ACM SIGART Bulletin} 2(2):7, 1991.
\bibitem{gjn} I.~P.~Gent, C.~Jefferson, and P.~Nightingale, ``Complexity of $n$-queens completion,'' \emph{Journal of Artificial Intelligence Research} 59:815--848, 2017.
\bibitem{hlm} E.~J.~Hoffman, J.~C.~Loessi, and R.~C.~Moore, ``Constructions for the solution of the $m$ queens problem,'' \emph{Mathematics Magazine} 42(2):66--72, 1969.
\bibitem{pauls} E.~Pauls, ``Das Maximalproblem der Damen auf dem Schachbrete,'' \emph{Deutsche Schachzeitung} 29:129--134, 257--267, 1874.
\bibitem{simkin} M.~Simkin, ``The number of $n$-queens configurations,'' arXiv:2107.13460, 2021.
\bibitem{mathlib} The Mathlib Community, \emph{Mathlib}, commit \texttt{a719ba5c31} (Lean v4.10.0), 2024.
\bibitem{lean} L.~de Moura and S.~Ullrich, ``The Lean 4 theorem prover and programming language,'' CADE 28, 2021.
\end{thebibliography}

\appendix

\section{Complete Lean source}
\label{app:lean}
\texttt{proofs/lean-project/ProofLab/NQueensTheorem.lean}, as integrated on the close-out
branch (identical to the reviewed commit \texttt{8ae158a}). Included verbatim from the
repository at build time.

\VerbatimInput[fontsize=\footnotesize,numbers=left,numbersep=4pt,frame=lines]{../../lean-project/ProofLab/NQueensTheorem.lean}

\section{Verification transcripts}
\label{app:transcripts}

\subsection{Reviewer's Lean checks}
Output of \texttt{lake env lean review\_checks.lean} (\texttt{review\_checks.out}):
\VerbatimInput[fontsize=\footnotesize]{../../../attacks/n-queens-20260912-180329/review_checks.out}

\subsection{Attack Lead's brute-force evidence (not proof)}
Output of \texttt{verify\_construction.py 1000} (\texttt{verify\_output.txt}). Note that this
script used signed integers and therefore could not detect the $\N$-truncation described in
Section~\ref{sec:erratum}.
\VerbatimInput[fontsize=\footnotesize]{../../../attacks/n-queens-20260912-180329/verify_output.txt}

\section{Ticket timeline}
\label{app:timeline}

\begin{center}\small
\begin{tabular}{@{}llll@{}}
\toprule
UTC & ticket & role & event \\
\midrule
17:37 & OPE-2 & board & names the problem; approves tooling \\
17:44 & OPE-2 & Director & Lean 4.10.0 installed and green; plan and child chain posted \\
17:51 & OPE-3 & Scout & dossier, pin, novelty screen (PR \#3) \\
17:53 & OPE-2 & Director & Gate 1: dossier approved; names frozen \\
17:58 & OPE-4 & Formalist & statement stub compiles with \lean{sorryAx} (PR \#4) \\
18:00 & OPE-2 & Director & Gate 2: stub approved; Attack Lead released \\
18:08 & OPE-5 & Attack Lead & construction, brute check, lemma list (PR \#5) \\
18:12 & OPE-2 & Director & Gate 4: plan approved; Formalist proof released \\
18:20 & OPE-6 & Formalist & both theorems closed, zero \lean{sorry} (PR \#6) \\
18:23 & OPE-2 & Director & Gate 6: rebuild and axioms re-checked; Reviewer released \\
18:28 & OPE-7 & Reviewer & approve with residuals (PR \#7) \\
18:34 & OPE-2 & Director & close-out: integration, erratum, ledger, claim packet (PR \#8) \\
\bottomrule
\end{tabular}
\end{center}

\end{document}
